\section*{الفصل \ref{ch:g4:what-plants-need} --- ما يحتاج إليه النبات}

\begin{solution}{exo:g4:what-plants-need:1}
الماء، والضوء، والهواء، والدفء، \omterm{prop:g4:what-plants-need:needs}{والأملاح المعدنية} (تُشرب مذابة في
ماء التربة).
\end{solution}

\begin{solution}{exo:g4:what-plants-need:2}
شيء واحد بالضبط --- وهو الشيء المختبَر؛ وكل ما عداه يجب أن يكون
متطابقًا. والأصيص الذي لم يُمسّ هو الشاهد.
\end{solution}

\begin{solution}{exo:g4:what-plants-need:3}
صفراء شاحبة بدل أن تكون خضراء، وطويلة ممتدة على نحو غريب، وضعيفة
--- ثم تستسلم. فقد كان الضوء ناقصًا، فلم تستطع \omterm{def:g1:plants-around-us:plant}{النبتة} أن تصنع
مادتها الخضراء ولا غذاءها.
\end{solution}

\begin{solution}{exo:g4:what-plants-need:4}
مقادير ضئيلة من مواد مغذية موجودة في التربة. وتشربها \omterm{def:g1:plants-around-us:parts}{الجذور} مذابة
في ماء التربة.
\end{solution}

\begin{solution}{exo:g4:what-plants-need:5}
لأن المحاصيل تحمل أملاح التربة المعدنية بعيدًا سنة بعد سنة؛ والسماد
يعيد تموينها. فهو يغذّي \emph{التربة}، ثم يحمل ماء التربة الأملاح
إلى \omterm{def:g1:plants-around-us:parts}{الجذور}.
\end{solution}

\begin{solution}{exo:g4:what-plants-need:6}
\omterm{def:g1:animals-around-us:animal}{الحيوان} ينمو من الكائنات الحية التي يأكلها؛ أما \omterm{def:g1:plants-around-us:plant}{النبات} فلا يأكل
أحدًا --- بل يصنع مادته بنفسه من الماء \omterm{prop:g4:what-plants-need:needs}{والأملاح المعدنية} والهواء
\omterm{prop:g3:balanced-diet:jobs}{وطاقة} الضوء التي تلتقطها أوراقه الخضراء.
\end{solution}

\begin{solution}{exo:g4:what-plants-need:7}
لا من التربة: بل من الماء ومن الهواء، والضوء يوفّر \omterm{prop:g3:balanced-diet:jobs}{الطاقة} ---
وبُنيت في \omterm{def:g1:plants-around-us:parts}{الأوراق} الخضراء. (ولم تسهم التربة إلا بالماء وبمقادير
صغيرة من \omterm{prop:g4:what-plants-need:needs}{الأملاح المعدنية}.)
\end{solution}

\begin{solution}{exo:g4:what-plants-need:8}
الجواب مفتوح. والمتوقع: الأصيص المسقيّ أخضر منتصب، وتوأمه الجاف
ذابل في أيام --- ولما كان الماء وحده هو المختلف، نُسب الفضل إلى
الماء.
\end{solution}

\begin{solution}{exo:g4:what-plants-need:9}
لا شاهد هناك: أصيص واحد، وكل شيء دفعةً واحدة --- فلعل الأصيص كان
ذا ضوء جيد وسقي دؤوب فحسب (فالمغنّون يزورون نباتاتهم كثيرًا!).
والنسخة المنصفة: أصيصان متطابقان، \omterm{def:g2:from-seed-to-plant:seed}{بالبذور} والتربة والنافذة والسقي
نفسها؛ ويُغنّى لأحدهما فحسب؛ ثم يقارَن بعد شهر. وعندئذ فقط يدل
الفرق على الغناء.
\end{solution}

\begin{solution}{exo:g4:what-plants-need:10}
الضوء: فهي تنمو مغمورة لكن في ماء مضاء، وأوراقها الخضراء ما زالت
تلتقط الضوء عبره. والماء حولها يوفّر أيضًا ما يوفّره ماء التربة ---
أملاحًا معدنية مذابة، تُشرب بسطحها كله بدل \omterm{def:g1:plants-around-us:parts}{الجذور}.
\end{solution}

\begin{solution}{exo:g4:what-plants-need:11}
\omterm{prop:g4:what-plants-need:needs}{الأملاح المعدنية}: فالماء والضوء والهواء والدفء متساوية، والرمل
النقي مع ماء المطر لا يوفّر أملاحًا تقريبًا. وقد نمت في البداية على
المدخرات المحزومة في أنسجتها هي (كما نمت \omterm{prop:g2:from-seed-to-plant:cycle}{البادرة} المدفونة على مخزن
\omterm{def:g2:from-seed-to-plant:seed}{البذرة}) --- فلما نفدت المدخرات ظهر النقص.
\end{solution}
